\chapter{Offline Reinforcement Learning}
\label{ch:offline_rl}
\noindent\textit{Lecture-aligned study notes: CS 8803 Lecture 11, Offline RL and (Deep) Offline RL (Fall 2026).}

\section{Problem Setup and the Main Difficulty}
Offline reinforcement learning seeks to learn a policy using only a fixed dataset
\[
\D=\{(s_i,a_i,r_i,s'_i,d_i)\}_{i=1}^{N},
\]
typically collected by a behavior policy $\pi_\beta$ or a mixture of behavior policies. During offline training, no new transitions can be requested. This is not synonymous with \emph{off-policy}: DQN and SAC can reuse data from a different policy while still collecting new transitions during online learning.

Behavior cloning (BC) maximizes $\E_{\D}[\log\pi_\theta(a\mid s)]$ without using rewards or Bellman backups. Offline RL additionally exploits observed rewards and long-horizon value estimation. It may \emph{stitch} useful portions of different trajectories where the state coverage permits, but it cannot reliably invent actions absent from the dataset.

\paragraph{Distribution shift and bootstrapping.}
A learned critic may assign arbitrarily optimistic values to out-of-distribution (OOD) state-action pairs. A greedy policy then exploits these errors:
\[
y=r+\gamma(1-d)\max_{a'}Q(s',a').
\]
If an unsupported action has estimated value $100$ while supported actions have values $5$ and $8$, the maximum selects $100$ even if its true value is small. Bellman backups propagate this error to earlier transitions. An online agent may correct the estimate with fresh experience; an offline learner cannot. Thus \emph{policy improvement must respect data support, or value estimation must be conservative enough to avoid exploiting uncertain extrapolations}.

\begin{pitfall}
Off-policy does not imply offline safety. Replay buffers and target networks stabilize learning but do not, by themselves, solve the OOD-action problem. An offline algorithm may also suffer state-distribution shift, approximation error, and insufficient coverage of promising actions.
\end{pitfall}

\section{Off-Policy Evaluation (OPE)}
OPE asks for the return $J(\pi)$ of a target policy $\pi$ using trajectories generated by behavior policy $\pi_\beta$, without executing $\pi$.

\subsection{Trajectory Importance Sampling}
Under appropriate support and trajectory-distribution assumptions, the importance ratio is
\[
\rho(\tau)=\prod_{t=0}^{T-1}
\frac{\pi(a_t\mid s_t)}{\pi_\beta(a_t\mid s_t)},
\qquad
\widehat J_{\rm IS}=\frac1N\sum_{i=1}^N \rho(\tau_i)R(\tau_i).
\]
The environment transition probabilities cancel in the likelihood ratio when target and behavior policies share the same environment and initial-state distribution. Support is essential: if $\pi(a\mid s)>0$ but $\pi_\beta(a\mid s)=0$ on a relevant state, the corresponding ratio cannot be identified from behavior data. Ratios multiply across time; for example, a ratio of $2$ in each of five steps yields $32$. This can lead to extreme variance and low effective sample size.

\subsection{Doubly Robust Evaluation}
In the contextual-bandit setting a standard doubly robust form is
\[
\widehat V_{\rm DR}(s,a,r)
=\widehat V(s)+
\frac{\pi(a\mid s)}{\pi_\beta(a\mid s)}
\bigl(r-\widehat Q(s,a)\bigr),
\quad
\widehat V(s)=\E_{a'\sim\pi(\cdot\mid s)}\widehat Q(s,a').
\]
The first term predicts return using a model; the second corrects its observed residual using importance weighting. Under suitable identification and estimation assumptions, correctness of either the propensity model or the outcome model can suffice for consistency. This does not guarantee low finite-sample variance. Sequential variants recursively combine value baselines and importance-weighted temporal corrections.

\subsection{Marginalized Importance Sampling}
Instead of multiplying policy ratios along entire trajectories, marginalized methods work with occupancy measures such as
\[
w(s,a)=\frac{d^\pi(s,a)}{d^{\pi_\beta}(s,a)}.
\]
Here $d^\pi$ is a normalized discounted occupancy measure when using discounted return conventions. The density ratio is estimated from distributional flow or stationarity relationships, as in the lecture's discussion of GenDICE. This can mitigate trajectory-product variance but estimating high-dimensional occupancy ratios remains difficult and requires coverage.

\section{Linear Offline Policy Evaluation: TD, LSTD, and LSTDQ}
\subsection{From a TD Target to an LSTD System}
For a fixed policy, approximate
\[
V_w(s)=\phi(s)^\top w,\qquad
\delta_i(w)=r_i+\gamma(1-d_i)V_w(s_i')-V_w(s_i).
\]
The one-step TD target is $y_i=r_i+\gamma(1-d_i)V_w(s_i')$ and the TD error is $\delta_i=y_i-V_w(s_i)$. For the exact value function the Bellman condition is
$\E[\delta\mid s]=0$, \emph{not} $\delta_i=0$ for every stochastic transition.

Linear semi-gradient TD(0) updates
\[
w\leftarrow w+\eta\,\delta_i(w)\phi(s_i),
\]
since $\nabla_w V_w(s_i)=\phi(s_i)$ and the bootstrap target is treated as fixed during the individual parameter update. LSTD instead solves the empirical feature-weighted equilibrium
\[
\sum_i\phi(s_i)\delta_i(w)=0,
\]
leading to the linear system
\[
\boxed{Aw=b},\qquad
A=\sum_i\phi(s_i)\bigl[\phi(s_i)-\gamma(1-d_i)\phi(s_i')\bigr]^\top,\quad
b=\sum_i\phi(s_i)r_i.
\]
Solve $w=A^{-1}b$ only if $A$ is nonsingular; in practice a linear solver or regularization is preferable. Multiplication by $\phi(s)$ enforces one balance equation per feature, not a separate zero-residual equation per transition.

\paragraph{Two-transition worked example.}
Use scalar features with $\gamma=0.5$ and two nonterminal samples:
\[
(\phi(s_1),\phi(s_1'),r_1)=(1,0,1),\qquad
(\phi(s_2),\phi(s_2'),r_2)=(2,1,2).
\]
Then
\[
A=1(1-0)+2(2-0.5)=4,\qquad b=1(1)+2(2)=5,
\quad w=1.25.
\]
The individual TD errors are $\delta_1=1-1.25=-0.25$ and
$\delta_2=2+0.5(1.25)-2(1.25)=0.125$. Neither is zero, but
$1(-0.25)+2(0.125)=0$.

\subsection{Projected Bellman Equation}
Let $T^\pi$ denote the policy Bellman operator. The exact fixed point satisfies $V^\pi=T^\pi V^\pi$. A linear feature space generally cannot represent this function exactly. Standard on-policy linear LSTD seeks the projected fixed point
\[
\boxed{V_w=\Pi T^\pi V_w},
\]
where $\Pi$ is projection onto the feature span under the relevant state weighting. This implies a feature-weighted Bellman-residual orthogonality condition. It differs from directly minimizing $\sum_i\delta_i(w)^2$: the full gradient of that objective differentiates through the next-state value, whereas semi-gradient TD and LSTD use current-state features in their normal equations. For off-policy or insufficiently covered data, additional assumptions and stability issues arise.

\subsection{LSTDQ and Least-Squares Policy Iteration (LSPI)}
For action values use $Q_w(s,a)=\phi(s,a)^\top w$. To evaluate a fixed policy $\pi$, set
\[
A_\pi=\sum_i\phi(s_i,a_i)
\left[\phi(s_i,a_i)-\gamma(1-d_i)\phi(s_i',\pi(s_i'))\right]^\top,\quad
b=\sum_i\phi(s_i,a_i)r_i.
\]
This is LSTDQ. LSPI alternates \emph{policy evaluation} using LSTDQ and \emph{greedy policy improvement}:
\[
w_k\leftarrow \operatorname{LSTDQ}(\D,\pi_k),\qquad
\pi_{k+1}(s)=\arg\max_a\phi(s,a)^\top w_k.
\]
The dataset stays fixed, but $A_{\pi_k}$ changes because the bootstrap action $\pi_k(s_i')$ changes. Greedy improvement can select unsupported actions and thereby reintroduce distribution shift.

\section{Deep Offline RL: Restrict the Policy or Correct the Critic}
\subsection{Explicit Policy Constraints and BEAR}
A generic constrained improvement step is
\[
\max_\pi \E_{s\sim\D,\ a\sim\pi(\cdot\mid s)}Q(s,a)
\quad\text{subject to}\quad
D(\pi(\cdot\mid s),\pi_\beta(\cdot\mid s))\leq\epsilon.
\]
BEAR (Bootstrapping Error Accumulation Reduction) emphasizes staying near the action \emph{support} of the behavior distribution, using a maximum mean discrepancy (MMD) based approximation. Matching support is not the same as matching all behavior probabilities: a new policy may put more mass on better supported actions. MMD is a surrogate, not a strict per-action support certificate. Behavior-model error and overly weak or strong constraints can respectively permit extrapolation or prevent improvement.

\subsection{Implicit Constraints: AWR and Advantage-Weighted BC}
KL-regularized improvement admits an exponential-reweighting form
\[
\pi^*(a\mid s)\propto\pi_\beta(a\mid s)
\exp\!\left(\frac{A(s,a)}{\lambda}\right).
\]
This motivates advantage-weighted regression (AWR):
\[
\boxed{\mathcal L_{\rm AWR}(\theta)=
-\E_{(s,a)\sim\D}
\left[\exp\!\left(\frac{\widehat A(s,a)}{\lambda}\right)
\log\pi_\theta(a\mid s)\right].}
\]
The actor learns via weighted log likelihood of \emph{dataset actions}, rather than directly differentiating $-Q(s,\pi_\theta(s))$. High-advantage examples have stronger influence. The critic estimation step may still query policy-generated OOD actions in its Bellman target, however. This limitation motivates IQL.

\subsection{Implicit Q-Learning (IQL)}
IQL learns a value function, a Q-function, and a policy. Its key idea is to perform critic backups without sampling a new action from the learned actor at $s'$.

\paragraph{Value loss: asymmetric expectile regression.}
Define the residual $u=Q_{\bar\phi}(s,a)-V_\psi(s)$ and
\[
\ell_2^\tau(u)=
\begin{cases}
\tau u^2,&u\geq0,\\
(1-\tau)u^2,&u<0.
\end{cases}
\qquad
\boxed{\mathcal L_V(\psi)=\E_{\D}[\ell_2^\tau(Q_{\bar\phi}(s,a)-V_\psi(s))].}
\]
For $\tau>0.5$, underestimating high Q-values is penalized more than overestimating low Q-values, so $V_\psi(s)$ moves toward the upper expectile of \emph{in-dataset} action values. With equally weighted Q-values $2$ and $8$ and $\tau=0.8$, the optimum between them solves
$0.2(V-2)+0.8(V-8)=0$, giving $V=6.8$ rather than the mean $5$ or maximum $8$.

\paragraph{Q loss: value-based backup.}
\[
\boxed{\mathcal L_Q(\phi)=
\E_{\D}\left[
\left(Q_\phi(s,a)-
\bigl[r+\gamma(1-d)V_\psi(s')\bigr]\right)^2
\right].}
\]
The target contains $V(s')$ rather than $\max_{a'}Q(s',a')$ or an expectation over actions sampled from the new policy. It does not directly query an OOD action value in the bootstrap step.

\paragraph{Policy loss: advantage-weighted extraction.}
\[
\boxed{\mathcal L_\pi(\theta)=
-\E_{\D}\left[
\exp\bigl(\beta[Q_\phi(s,a)-V_\psi(s)]\bigr)
\log\pi_\theta(a\mid s)\right].}
\]
The exponential weight is treated as constant during the actor update. $\tau$ controls how optimistic $V$ is within the observed-action distribution; $\beta$ controls the strength of advantage preference in the actor. They are distinct knobs. In implementations, advantage-weight clipping, target Q-networks and double critics are common stability details. IQL cannot reliably discover wholly unsupported high-value actions, and errors in in-dataset Q estimates still matter.

\subsection{Conservative Q-Learning (CQL)}
Rather than avoiding OOD Q queries, CQL makes critic values conservative. For a discrete action space, a common form is
\[
\boxed{
\mathcal L_{\rm CQL}=
\mathcal L_{\rm TD}
+\alpha\left[
\E_{s\sim\D}\log\sum_a\exp Q(s,a)
-\E_{(s,a)\sim\D}Q(s,a)
\right].
}
\]
The log-sum-exp term pressures large action values downward, while the data-action term relatively preserves observed actions. This discourages a learned policy from exploiting spuriously high Q-values. Continuous-action variants approximate the action expectation using samples, rather than enumerate every action. Too little regularization fails to address overestimation; too much can suppress genuinely useful actions. CQL does not guarantee every OOD action's Q is below every in-dataset action's Q.

\section{Comparison, Training Pitfalls, and Extensions}
\begin{center}\small
\begin{tabularx}{\textwidth}{@{}p{0.15\textwidth}X X@{}}\toprule
Method & Main intervention & Important tradeoff or failure mode\\\midrule
BEAR & Restrict actor near behavior action support (MMD surrogate) & Coverage estimation; too strong a constraint reduces improvement\\
AWR & Weighted behavior cloning from advantages & Critic backups may still evaluate policy-generated OOD actions\\
IQL & Expectile $V(s)$, $r+\gamma V(s')$ backup, weighted BC & Quality of dataset actions and Q estimates; expectile and weight tuning\\
CQL & Conservative penalty on learned Q-values & Excessive pessimism and costly action sampling\\\bottomrule
\end{tabularx}
\end{center}

\paragraph{Model-based and offline-to-online extensions.}
The lectures also discuss COMBO, which combines learned dynamics and conservative offline value learning, and offline pretraining followed by online fine-tuning (including AWAC/IQL). Model-generated rollouts provide extra coverage but introduce model bias; online fine-tuning changes the setting by allowing new interaction.

\paragraph{Training checklist.}
Before choosing a method, check behavior support and dataset quality; distinguish action support from distribution matching; monitor extrapolation and critic divergence; tune $\tau$, advantage-weight temperature, MMD strength or CQL $\alpha$ with care; and do not assume that a lower training TD loss implies reliable evaluation on unseen actions.

\section{Review Questions}
\begin{enumerate}
\item Explain why off-policy learning and offline RL are not equivalent.
\item Show how OOD overestimation enters a max-Q Bellman target and propagates through time.
\item Why can trajectory importance sampling have high variance? State the support requirement.
\item Explain the correction term in a doubly robust estimator and how occupancy ratios differ from trajectory products.
\item Derive $Aw=b$ from $\sum_i\phi(s_i)\delta_i=0$ and explain why each $\delta_i$ need not vanish.
\item Why can LSPI change its LSTDQ system without changing the dataset?
\item Contrast BEAR's support constraint, IQL's expectile backup, and CQL's critic penalty.
\item For Q-values $2$ and $8$, why does expectile $\tau=0.8$ produce $V=6.8$?
\item Which of IQL's three losses uses dataset actions in policy extraction, and what differentiates it from SAC's actor update?
\end{enumerate}

\paragraph{Source and scope.}
Primary sources: CS 8803 Fall 2026 Lecture 11 slides, \emph{Offline RL} and \emph{(Deep) Offline RL}. Lecture-discussed algorithms include LSTD/LSTDQ/LSPI, BEAR, advantage-weighted policy learning, IQL, CQL, COMBO, and OPE approaches such as importance sampling and doubly robust estimation. The worked examples and explanations are instructional syntheses, not reproduced paper results.
